% Shared by main.tex and main_synthesis.tex; campaign audited 22 September 2026.
% Figure generation: scripts/reports/make_closed_loop_stress_figures.py.
\subsection{Closed-loop and cross-directional stress tests: preliminary results}
\label{sec:adr-closed-loop-stress}

This extension tests three distinct transport/diffusion arrangements on a
more tightly pinched annulus. The completed strong-preset campaign of
22 September 2026 is reported separately from the preceding solver rankings:
eight of 60 comparisons passed, all on the orthogonal control. The figures
below show analytic input data and a saved campaign mesh, not a computed
HDG solution.

\paragraph{Common geometry and exact solution.}
For the main severity level, let
\begin{equation}
 \begin{gathered}
 R(\phi)=1+0.35\cos(9\phi),\qquad
 \Omega=\{(r,\phi):0.63<r<R(\phi)\},\\
 \rho=\frac{r-0.63}{R(\phi)-0.63},\qquad
 \chi=\rho+0.05\sin(5\phi)\sin(2\pi\rho),\\
 u_\star=\sin(2\pi\chi)
 +\tfrac14\bigl[\sin(6\pi x)\sin(5\pi y)
 +0.35\sin(11\pi x)\sin(9\pi y)\bigr].
 \end{gathered}
 \label{eq:adr-closed-loop-exact}
\end{equation}
There are nine nominal necks of width $0.02$, compared with $0.07$ in the
earlier five-lobed geometry. All variants solve
\begin{equation}
 \begin{gathered}
 -\nabla\!\cdot(K\nabla u)+\nabla\!\cdot(\boldsymbol\beta u)+c u=f,
 \qquad u|_{\partial\Omega}=u_\star,\\
 \epsilon=10^{-6},\quad U=50,\quad c=10^{-3}.
 \end{gathered}
 \label{eq:adr-closed-loop-pde}
\end{equation}
The divergence-free velocities give the manufactured source
$f=-K:\nabla^2u_\star-(\nabla\!\cdot K)\cdot\nabla u_\star
+\boldsymbol\beta\cdot\nabla u_\star+c u_\star$;
the derivative-of-$K$ term is retained for the curved tensors.

\begin{figure}[htbp]
 \centering
 \includegraphics[width=\linewidth]{\ADRResultsFigurePath/closed_loop_stress/geometry_fields.pdf}
 \caption{Main-level stress geometry, common to all three variants:
 (a) the nine-lobed annulus, (b) a magnified minimum-width neck, and
 (c) the same neck in the saved $98\,699$-triangle campaign mesh.
 The first two panels show ideal smooth walls, not a computational mesh;
 the third preserves the actual affine triangles, with nonuniform sizes
 and shapes, at equal axis scale. No artificial deformation or remeshing
 is applied; a mesh close-up alone does not establish spatial resolution.}
 \label{fig:adr-closed-loop-fields}
\end{figure}

\paragraph{Three transport arrangements.}
For \emph{trapping} and \emph{crossing}, define
$b=\nabla^\perp\chi/|\nabla\chi|$,
$K=\epsilon I+(1-\epsilon)bb^T$ and
$\boldsymbol\beta=(U/C_\psi)\nabla^\perp\psi$, where
$\nabla^\perp=(\partial_y,-\partial_x)$ and $C_\psi$ is a
mesh-independent, empirically converged estimate of
$\|\nabla\psi\|_{L^\infty(\Omega)}$, frozen across the mesh/order sweep.
The streamfunctions are
\begin{align}
 \psi_{\rm trap}&=-\frac{\cos(8\pi\chi)}{8\pi},\label{eq:adr-closed-loop-trap}\\
 \psi_{\rm cross}&=16\rho^2(1-\rho)^2\left[
 \frac{\sin(13\pi x)\sin(11\pi y)}{13\pi}
 +0.35\frac{\sin(17\pi x)\sin(19\pi y)}{19\pi}\right].
 \label{eq:adr-closed-loop-cross}
\end{align}
Trapping follows the closed diffusion paths; crossing circulation generally
cuts across them. The \emph{orthogonal} control instead fixes
$e_d=(\cos(\pi/7),\sin(\pi/7))^T$ and
$e_t=(-\sin(\pi/7),\cos(\pi/7))^T$, with
\begin{equation}
 K_\perp=e_de_d^T+\epsilon e_te_t^T,\qquad
 \boldsymbol\beta_\perp=
 \frac{U}{1.8}\bigl[1+0.8\sin(13\pi e_d\cdot(x,y))\bigr]e_t.
 \label{eq:adr-closed-loop-orthogonal}
\end{equation}
This is nonvanishing through-flow exactly perpendicular to strong diffusion,
not a recirculating case. Exact Dirichlet data are imposed on both polygonal
walls in every variant; analytic wall tangency in the first two does not
imply exact impermeability of their affine boundary approximations.

\paragraph{Distinguishing the three operators.}
The manufactured $u_\star$ is identical in all three cases, so plotting it
three times cannot distinguish the tests. Instead,
Figure~\ref{fig:adr-closed-loop-operator-terms} shows the signed PDE
contributions on that common field,
\[
 A_\star=\boldsymbol\beta\cdot\nabla u_\star,\qquad
 D_\star=-K:\nabla^2u_\star-(\nabla\cdot K)\cdot\nabla u_\star.
\]
The divergence-free velocities make $A_\star=\nabla\cdot(\boldsymbol\beta u_\star)$.
The tensor-derivative term is retained in $D_\star$.
Arrows follow the trapping loops, the crossing circulation, or the
orthogonal through-flow; headless segments identify the strong-diffusion
axis. Trapping and crossing share $K$ and therefore the same $D_\star$,
whereas the orthogonal control has a fixed diffusion axis.
These views explain the operator differences that the solver comparison
tests; local term amplitudes alone are not condition-number estimates.

\begin{figure}[htbp]
 \centering
 \includegraphics[width=\linewidth]{\ADRResultsFigurePath/closed_loop_stress/operator_terms.pdf}
 \caption{Analytic operator contributions on the common exact solution:
 columns are trapping, crossing and orthogonal; the top row shows advection
 $A_\star$ and the bottom row diffusion $D_\star$. Colors are signed, with
 one shared symmetric-log scale per row (linear for $|A_\star|,|D_\star|\leq1$);
 the two rows have different color limits. Equal-length arrows show only
 advection direction, not speed; headless segments show the strong-diffusion
 direction, not diffusive flux. Frozen campaign velocity normalizations
 are used. The matching trapping/crossing diffusion panels are intentional.
 These are analytic input diagnostics, not computed solutions or residuals.}
 \label{fig:adr-closed-loop-operator-terms}
\end{figure}

\paragraph{Why these are demanding.}
The tensor has anisotropy ratio $10^6$, reaction is weak, and the narrow
necks compress geometric and coefficient scales. In the trapping case,
both strong diffusion and advection annihilate functions $g(\chi)$ supported
away from the walls, leaving
$\mathcal{L}g(\chi)=-\epsilon\Delta g(\chi)+c\,g(\chi)$:
cross-loop modes are only weakly controlled. Crossing cells and perpendicular
through-flow instead test nonsymmetric couplings that a diffusion-based
correction may not capture; neither a difficulty ordering nor a solver
failure follows in advance, and oscillatory forcing alone does not alter
matrix conditioning.

\input{\ADRResultsPath/closed_loop_stress/results.tex}
